% =====================================================================
%  Anatomía de un LLM — Fase 0 — Módulo 0.8
%  Combinatoria y probabilidad
%  Compilar: latexmk -r ../latexmkrc P08-probabilidad.tex
% =====================================================================
\documentclass[11pt,a4paper,twoside,openany]{book}
\usepackage{../estilo-apuntes}

\newcommand{\ApuntesCurso}{Anatomía de un LLM}
\newcommand{\ApuntesModulo}{Módulo 0.8}

% Anclas de hiperenlace únicas por tema para los enunciados
\makeatletter
\@for\nombre:=definicion,teorema,proposicion,lema,corolario,ejemplo,observacion\do{%
  \expandafter\xdef\csname theH\nombre\endcsname{\noexpand\thechapter.\noexpand\arabic{definicion}}}
\makeatother

% --- Notación propia del módulo -------------------------------------------
\newcommand{\N}{\mathbb{N}}
\newcommand{\PP}{\mathbb{P}}
\newcommand{\Nor}{\mathcal{N}}
\newcommand{\Xbar}{\overline{X}}
\newcommand{\dsim}{\mathbin{\triangle}}
\DeclareMathOperator{\Var}{Var}
\DeclareMathOperator{\Ber}{Ber}
\DeclareMathOperator{\Bin}{Bin}
\DeclareMathOperator{\Geom}{Geom}
\DeclareMathOperator{\Poi}{Poisson}
\DeclareMathOperator{\Cat}{Cat}
\DeclareMathOperator{\Unif}{Unif}
\DeclareMathOperator{\Expo}{Exp}

\hypersetup{pdftitle={Módulo 0.8. Combinatoria y probabilidad (resumen)},pdfauthor={Anatomía de un LLM}}

\makeatletter
\newlabel{ej:t1-contrasenas}{{1.5}{2}{Contraseñas}{ejemplo.1.5}{}}
\newlabel{ej:t1-contrasenas@cref}{{[ejemplo][5][]1.5}{[1][2][]2}}
\newlabel{ej:t1-bits}{{1.7}{2}{Subconjuntos y cadenas de bits}{ejemplo.1.7}{}}
\newlabel{ej:t1-bits@cref}{{[ejemplo][7][]1.7}{[1][2][]2}}
\newlabel{ej:t1-divisores}{{1.10}{3}{}{ejemplo.1.10}{}}
\newlabel{ej:t1-divisores@cref}{{[ejemplo][10][]1.10}{[1][3][]3}}
\newlabel{ej:t1-circular}{{1.13}{3}{Permutaciones circulares}{ejemplo.1.13}{}}
\newlabel{ej:t1-circular@cref}{{[ejemplo][13][]1.13}{[1][3][]3}}
\newlabel{eje:1-reglas}{{1.1}{4}{}{ejercicio.1.1}{}}
\newlabel{eje:1-reglas@cref}{{[ejercicio][1][1]1.1}{[1][4][]4}}
\newlabel{eje:1-vocabulario}{{1.2}{4}{}{ejercicio.1.2}{}}
\newlabel{eje:1-vocabulario@cref}{{[ejercicio][2][1]1.2}{[1][4][]4}}
\newlabel{eje:1-contrasenas}{{1.3}{4}{}{ejercicio.1.3}{}}
\newlabel{eje:1-contrasenas@cref}{{[ejercicio][3][1]1.3}{[1][4][]4}}
\newlabel{eje:1-divisibles}{{1.4}{4}{}{ejercicio.1.4}{}}
\newlabel{eje:1-divisibles@cref}{{[ejercicio][4][1]1.4}{[1][4][]4}}
\newlabel{eje:1-palindromos}{{1.5}{4}{}{ejercicio.1.5}{}}
\newlabel{eje:1-palindromos@cref}{{[ejercicio][5][1]1.5}{[1][4][]4}}
\newlabel{eje:1-mesa}{{1.6}{5}{}{ejercicio.1.6}{}}
\newlabel{eje:1-mesa@cref}{{[ejercicio][6][1]1.6}{[1][5][]5}}
\newlabel{eje:1-cifras}{{1.7}{5}{}{ejercicio.1.7}{}}
\newlabel{eje:1-cifras@cref}{{[ejercicio][7][1]1.7}{[1][5][]5}}
\newlabel{eje:1-funciones}{{1.8}{5}{}{ejercicio.1.8}{}}
\newlabel{eje:1-funciones@cref}{{[ejercicio][8][1]1.8}{[1][5][]5}}
\newlabel{eje:1-sombreros}{{1.9}{5}{}{ejercicio.1.9}{}}
\newlabel{eje:1-sombreros@cref}{{[ejercicio][9][1]1.9}{[1][5][]5}}
\newlabel{eje:1-encuesta}{{1.10}{5}{}{ejercicio.1.10}{}}
\newlabel{eje:1-encuesta@cref}{{[ejercicio][10][1]1.10}{[1][5][]5}}
\newlabel{eje:1-producto-ind}{{1.11}{5}{}{ejercicio.1.11}{}}
\newlabel{eje:1-producto-ind@cref}{{[ejercicio][11][1]1.11}{[1][5][]5}}
\newlabel{eje:1-fibonacci}{{1.12}{5}{}{ejercicio.1.12}{}}
\newlabel{eje:1-fibonacci@cref}{{[ejercicio][12][1]1.12}{[1][5][]5}}
\newlabel{eje:1-ie-cuatro}{{1.13}{5}{}{ejercicio.1.13}{}}
\newlabel{eje:1-ie-cuatro@cref}{{[ejercicio][13][1]1.13}{[1][5][]5}}
\newlabel{eje:1-collares}{{1.14}{5}{}{ejercicio.1.14}{}}
\newlabel{eje:1-collares@cref}{{[ejercicio][14][1]1.14}{[1][5][]5}}
\newlabel{ej:t2-caminos}{{2.12}{9}{Caminos en una cuadrícula}{ejemplo.2.12}{}}
\newlabel{ej:t2-caminos@cref}{{[ejemplo][12][]2.12}{[1][9][]9}}
\newlabel{ej:t2-desarreglos}{{2.15}{9}{Desarreglos}{ejemplo.2.15}{}}
\newlabel{ej:t2-desarreglos@cref}{{[ejemplo][15][]2.15}{[1][9][]9}}
\newlabel{eje:2-calcular}{{2.1}{10}{}{ejercicio.2.1}{}}
\newlabel{eje:2-calcular@cref}{{[ejercicio][1][2]2.1}{[1][10][]10}}
\newlabel{eje:2-modelo}{{2.2}{10}{}{ejercicio.2.2}{}}
\newlabel{eje:2-modelo@cref}{{[ejercicio][2][2]2.2}{[1][10][]10}}
\newlabel{eje:2-poker}{{2.3}{10}{}{ejercicio.2.3}{}}
\newlabel{eje:2-poker@cref}{{[ejercicio][3][2]2.3}{[1][10][]10}}
\newlabel{eje:2-equipos}{{2.4}{10}{}{ejercicio.2.4}{}}
\newlabel{eje:2-equipos@cref}{{[ejercicio][4][2]2.4}{[1][10][]10}}
\newlabel{eje:2-anagramas}{{2.5}{11}{}{ejercicio.2.5}{}}
\newlabel{eje:2-anagramas@cref}{{[ejercicio][5][2]2.5}{[1][10][]11}}
\newlabel{eje:2-coeficientes}{{2.6}{11}{}{ejercicio.2.6}{}}
\newlabel{eje:2-coeficientes@cref}{{[ejercicio][6][2]2.6}{[1][11][]11}}
\newlabel{eje:2-ecuaciones}{{2.7}{11}{}{ejercicio.2.7}{}}
\newlabel{eje:2-ecuaciones@cref}{{[ejercicio][7][2]2.7}{[1][11][]11}}
\newlabel{eje:2-caminos}{{2.8}{11}{}{ejercicio.2.8}{}}
\newlabel{eje:2-caminos@cref}{{[ejercicio][8][2]2.8}{[1][11][]11}}
\newlabel{eje:2-estaciones}{{2.9}{11}{}{ejercicio.2.9}{}}
\newlabel{eje:2-estaciones@cref}{{[ejercicio][9][2]2.9}{[1][11][]11}}
\newlabel{eje:2-desarreglos}{{2.10}{11}{}{ejercicio.2.10}{}}
\newlabel{eje:2-desarreglos@cref}{{[ejercicio][10][2]2.10}{[1][11][]11}}
\newlabel{eje:2-ositos}{{2.11}{11}{}{ejercicio.2.11}{}}
\newlabel{eje:2-ositos@cref}{{[ejercicio][11][2]2.11}{[1][11][]11}}
\newlabel{eje:2-palo}{{2.12}{11}{}{ejercicio.2.12}{}}
\newlabel{eje:2-palo@cref}{{[ejercicio][12][2]2.12}{[1][11][]11}}
\newlabel{eje:2-sumak}{{2.13}{11}{}{ejercicio.2.13}{}}
\newlabel{eje:2-sumak@cref}{{[ejercicio][13][2]2.13}{[1][11][]11}}
\newlabel{eje:2-vandermonde}{{2.14}{11}{}{ejercicio.2.14}{}}
\newlabel{eje:2-vandermonde@cref}{{[ejercicio][14][2]2.14}{[1][11][]11}}
\newlabel{eje:2-pares}{{2.15}{12}{}{ejercicio.2.15}{}}
\newlabel{eje:2-pares@cref}{{[ejercicio][15][2]2.15}{[1][11][]12}}
\newlabel{ej:t3-omega}{{3.2}{13}{}{ejemplo.3.2}{}}
\newlabel{ej:t3-omega@cref}{{[ejemplo][2][]3.2}{[1][13][]13}}
\newlabel{ej:t3-ases}{{3.9}{15}{Ordenado o no ordenado}{ejemplo.3.9}{}}
\newlabel{ej:t3-ases@cref}{{[ejemplo][9][]3.9}{[1][15][]15}}
\newlabel{fig:t3-cumple}{{3.1}{16}{Probabilidad $p_n$ de que al menos dos de $n$ personas compartan cumpleaños ($N=365$)}{figure.3.1}{}}
\newlabel{fig:t3-cumple@cref}{{[figure][1][3]3.1}{[1][16][]16}}
\newlabel{eje:3-espacio}{{3.1}{17}{}{ejercicio.3.1}{}}
\newlabel{eje:3-espacio@cref}{{[ejercicio][1][3]3.1}{[1][16][]17}}
\newlabel{eje:3-sucesos}{{3.2}{17}{}{ejercicio.3.2}{}}
\newlabel{eje:3-sucesos@cref}{{[ejercicio][2][3]3.2}{[1][17][]17}}
\newlabel{eje:3-monedas}{{3.3}{17}{}{ejercicio.3.3}{}}
\newlabel{eje:3-monedas@cref}{{[ejercicio][3][3]3.3}{[1][17][]17}}
\newlabel{eje:3-dados}{{3.4}{17}{}{ejercicio.3.4}{}}
\newlabel{eje:3-dados@cref}{{[ejercicio][4][3]3.4}{[1][17][]17}}
\newlabel{eje:3-cuatro-dados}{{3.5}{17}{}{ejercicio.3.5}{}}
\newlabel{eje:3-cuatro-dados@cref}{{[ejercicio][5][3]3.5}{[1][17][]17}}
\newlabel{eje:3-horarios}{{3.6}{17}{}{ejercicio.3.6}{}}
\newlabel{eje:3-horarios@cref}{{[ejercicio][6][3]3.6}{[1][17][]17}}
\newlabel{eje:3-robos}{{3.7}{17}{}{ejercicio.3.7}{}}
\newlabel{eje:3-robos@cref}{{[ejercicio][7][3]3.7}{[1][17][]17}}
\newlabel{eje:3-seis-cartas}{{3.8}{17}{}{ejercicio.3.8}{}}
\newlabel{eje:3-seis-cartas@cref}{{[ejercicio][8][3]3.8}{[1][17][]17}}
\newlabel{eje:3-cumple}{{3.9}{17}{}{ejercicio.3.9}{}}
\newlabel{eje:3-cumple@cref}{{[ejercicio][9][3]3.9}{[1][17][]17}}
\newlabel{eje:3-urna}{{3.10}{18}{}{ejercicio.3.10}{}}
\newlabel{eje:3-urna@cref}{{[ejercicio][10][3]3.10}{[1][17][]18}}
\newlabel{eje:3-geometrica}{{3.11}{18}{}{ejercicio.3.11}{}}
\newlabel{eje:3-geometrica@cref}{{[ejercicio][11][3]3.11}{[1][18][]18}}
\newlabel{eje:3-bonferroni}{{3.12}{18}{}{ejercicio.3.12}{}}
\newlabel{eje:3-bonferroni@cref}{{[ejercicio][12][3]3.12}{[1][18][]18}}
\newlabel{eje:3-interseccion}{{3.13}{18}{}{ejercicio.3.13}{}}
\newlabel{eje:3-interseccion@cref}{{[ejercicio][13][3]3.13}{[1][18][]18}}
\newlabel{eje:3-distancia}{{3.14}{18}{}{ejercicio.3.14}{}}
\newlabel{eje:3-distancia@cref}{{[ejercicio][14][3]3.14}{[1][18][]18}}
\newlabel{eje:3-hash}{{3.15}{18}{}{ejercicio.3.15}{}}
\newlabel{eje:3-hash@cref}{{[ejercicio][15][3]3.15}{[1][18][]18}}
\newlabel{eje:3-laplace}{{3.16}{18}{}{ejercicio.3.16}{}}
\newlabel{eje:3-laplace@cref}{{[ejercicio][16][3]3.16}{[1][18][]18}}
\newlabel{ej:t4-hijos}{{4.3}{19}{}{ejemplo.4.3}{}}
\newlabel{ej:t4-hijos@cref}{{[ejemplo][3][]4.3}{[1][19][]19}}
\newlabel{ej:t4-diagnostico}{{4.9}{21}{Prueba diagnóstica}{ejemplo.4.9}{}}
\newlabel{ej:t4-diagnostico@cref}{{[ejemplo][9][]4.9}{[1][20][]21}}
\newlabel{ej:t4-monty}{{4.10}{21}{Monty Hall}{ejemplo.4.10}{}}
\newlabel{ej:t4-monty@cref}{{[ejemplo][10][]4.10}{[1][21][]21}}
\newlabel{ej:t4-dosados}{{4.13}{21}{Independencia dos a dos sin independencia mutua}{ejemplo.4.13}{}}
\newlabel{ej:t4-dosados@cref}{{[ejemplo][13][]4.13}{[1][21][]21}}
\newlabel{eje:4-tabla}{{4.1}{22}{}{ejercicio.4.1}{}}
\newlabel{eje:4-tabla@cref}{{[ejercicio][1][4]4.1}{[1][22][]22}}
\newlabel{eje:4-cachorros}{{4.2}{22}{}{ejercicio.4.2}{}}
\newlabel{eje:4-cachorros@cref}{{[ejercicio][2][4]4.2}{[1][22][]22}}
\newlabel{eje:4-spam}{{4.3}{22}{}{ejercicio.4.3}{}}
\newlabel{eje:4-spam@cref}{{[ejercicio][3][4]4.3}{[1][22][]22}}
\newlabel{eje:4-dos-pruebas}{{4.4}{23}{}{ejercicio.4.4}{}}
\newlabel{eje:4-dos-pruebas@cref}{{[ejercicio][4][4]4.4}{[1][22][]23}}
\newlabel{eje:4-piezas}{{4.5}{23}{}{ejercicio.4.5}{}}
\newlabel{eje:4-piezas@cref}{{[ejercicio][5][4]4.5}{[1][23][]23}}
\newlabel{eje:4-perro}{{4.6}{23}{}{ejercicio.4.6}{}}
\newlabel{eje:4-perro@cref}{{[ejercicio][6][4]4.6}{[1][23][]23}}
\newlabel{eje:4-monty7}{{4.7}{23}{}{ejercicio.4.7}{}}
\newlabel{eje:4-monty7@cref}{{[ejercicio][7][4]4.7}{[1][23][]23}}
\newlabel{eje:4-monty-p}{{4.8}{23}{}{ejercicio.4.8}{}}
\newlabel{eje:4-monty-p@cref}{{[ejercicio][8][4]4.8}{[1][23][]23}}
\newlabel{eje:4-paridad}{{4.9}{23}{}{ejercicio.4.9}{}}
\newlabel{eje:4-paridad@cref}{{[ejercicio][9][4]4.9}{[1][23][]23}}
\newlabel{eje:4-ajedrez}{{4.10}{23}{}{ejercicio.4.10}{}}
\newlabel{eje:4-ajedrez@cref}{{[ejercicio][10][4]4.10}{[1][23][]23}}
\newlabel{eje:4-si-mismo}{{4.11}{24}{}{ejercicio.4.11}{}}
\newlabel{eje:4-si-mismo@cref}{{[ejercicio][11][4]4.11}{[1][24][]24}}
\newlabel{eje:4-union-indep}{{4.12}{24}{}{ejercicio.4.12}{}}
\newlabel{eje:4-union-indep@cref}{{[ejercicio][12][4]4.12}{[1][24][]24}}
\newlabel{eje:4-producto-triple}{{4.13}{24}{}{ejercicio.4.13}{}}
\newlabel{eje:4-producto-triple@cref}{{[ejercicio][13][4]4.13}{[1][24][]24}}
\newlabel{eje:4-razones}{{4.14}{24}{}{ejercicio.4.14}{}}
\newlabel{eje:4-razones@cref}{{[ejercicio][14][4]4.14}{[1][24][]24}}
\newlabel{eje:4-simetria}{{4.15}{24}{}{ejercicio.4.15}{}}
\newlabel{eje:4-simetria@cref}{{[ejercicio][15][4]4.15}{[1][24][]24}}
\newlabel{ej:t5-suma}{{5.2}{25}{}{ejemplo.5.2}{}}
\newlabel{ej:t5-suma@cref}{{[ejemplo][2][]5.2}{[1][25][]25}}
\newlabel{ej:t5-fijos}{{5.8}{26}{Puntos fijos}{ejemplo.5.8}{}}
\newlabel{ej:t5-fijos@cref}{{[ejemplo][8][]5.8}{[1][26][]26}}
\newlabel{eje:5-masa}{{5.1}{29}{}{ejercicio.5.1}{}}
\newlabel{eje:5-masa@cref}{{[ejercicio][1][5]5.1}{[1][28][]29}}
\newlabel{eje:5-diferencia}{{5.2}{29}{}{ejercicio.5.2}{}}
\newlabel{eje:5-diferencia@cref}{{[ejercicio][2][5]5.2}{[1][29][]29}}
\newlabel{eje:5-binomial}{{5.3}{29}{}{ejercicio.5.3}{}}
\newlabel{eje:5-binomial@cref}{{[ejercicio][3][5]5.3}{[1][29][]29}}
\newlabel{eje:5-geometrica}{{5.4}{29}{}{ejercicio.5.4}{}}
\newlabel{eje:5-geometrica@cref}{{[ejercicio][4][5]5.4}{[1][29][]29}}
\newlabel{eje:5-poisson-aprox}{{5.5}{29}{}{ejercicio.5.5}{}}
\newlabel{eje:5-poisson-aprox@cref}{{[ejercicio][5][5]5.5}{[1][29][]29}}
\newlabel{eje:5-cumple-esp}{{5.6}{29}{}{ejercicio.5.6}{}}
\newlabel{eje:5-cumple-esp@cref}{{[ejercicio][6][5]5.6}{[1][29][]29}}
\newlabel{eje:5-hash}{{5.7}{29}{}{ejercicio.5.7}{}}
\newlabel{eje:5-hash@cref}{{[ejercicio][7][5]5.7}{[1][29][]29}}
\newlabel{eje:5-hijos}{{5.8}{29}{}{ejercicio.5.8}{}}
\newlabel{eje:5-hijos@cref}{{[ejercicio][8][5]5.8}{[1][29][]29}}
\newlabel{eje:5-inversiones}{{5.9}{29}{}{ejercicio.5.9}{}}
\newlabel{eje:5-inversiones@cref}{{[ejercicio][9][5]5.9}{[1][29][]29}}
\newlabel{eje:5-categorica}{{5.10}{29}{}{ejercicio.5.10}{}}
\newlabel{eje:5-categorica@cref}{{[ejercicio][10][5]5.10}{[1][29][]29}}
\newlabel{eje:5-passk}{{5.11}{30}{}{ejercicio.5.11}{}}
\newlabel{eje:5-passk@cref}{{[ejercicio][11][5]5.11}{[1][29][]30}}
\newlabel{eje:5-colas}{{5.12}{30}{}{ejercicio.5.12}{}}
\newlabel{eje:5-colas@cref}{{[ejercicio][12][5]5.12}{[1][30][]30}}
\newlabel{eje:5-poisson}{{5.13}{30}{}{ejercicio.5.13}{}}
\newlabel{eje:5-poisson@cref}{{[ejercicio][13][5]5.13}{[1][30][]30}}
\newlabel{eje:5-sin-memoria}{{5.14}{30}{}{ejercicio.5.14}{}}
\newlabel{eje:5-sin-memoria@cref}{{[ejercicio][14][5]5.14}{[1][30][]30}}
\newlabel{eje:5-geom-cond}{{5.15}{30}{}{ejercicio.5.15}{}}
\newlabel{eje:5-geom-cond@cref}{{[ejercicio][15][5]5.15}{[1][30][]30}}
\newlabel{ej:t6-regla}{{6.9}{33}{Regla 68--95--99,7}{ejemplo.6.9}{}}
\newlabel{ej:t6-regla@cref}{{[ejemplo][9][]6.9}{[1][32][]33}}
\newlabel{eje:6-conceptos}{{6.1}{34}{}{ejercicio.6.1}{}}
\newlabel{eje:6-conceptos@cref}{{[ejercicio][1][6]6.1}{[1][34][]34}}
\newlabel{eje:6-densidad}{{6.2}{34}{}{ejercicio.6.2}{}}
\newlabel{eje:6-densidad@cref}{{[ejercicio][2][6]6.2}{[1][34][]34}}
\newlabel{eje:6-pasas}{{6.3}{34}{}{ejercicio.6.3}{}}
\newlabel{eje:6-pasas@cref}{{[ejercicio][3][6]6.3}{[1][34][]34}}
\newlabel{eje:6-exponencial}{{6.4}{35}{}{ejercicio.6.4}{}}
\newlabel{eje:6-exponencial@cref}{{[ejercicio][4][6]6.4}{[1][34][]35}}
\newlabel{eje:6-autobus}{{6.5}{35}{}{ejercicio.6.5}{}}
\newlabel{eje:6-autobus@cref}{{[ejercicio][5][6]6.5}{[1][35][]35}}
\newlabel{eje:6-tipificar}{{6.6}{35}{}{ejercicio.6.6}{}}
\newlabel{eje:6-tipificar@cref}{{[ejercicio][6][6]6.6}{[1][35][]35}}
\newlabel{eje:6-gestacion}{{6.7}{35}{}{ejercicio.6.7}{}}
\newlabel{eje:6-gestacion@cref}{{[ejercicio][7][6]6.7}{[1][35][]35}}
\newlabel{eje:6-palo}{{6.8}{35}{}{ejercicio.6.8}{}}
\newlabel{eje:6-palo@cref}{{[ejercicio][8][6]6.8}{[1][35][]35}}
\newlabel{eje:6-markov}{{6.9}{35}{}{ejercicio.6.9}{}}
\newlabel{eje:6-markov@cref}{{[ejercicio][9][6]6.9}{[1][35][]35}}
\newlabel{eje:6-n-cheb}{{6.10}{35}{}{ejercicio.6.10}{}}
\newlabel{eje:6-n-cheb@cref}{{[ejercicio][10][6]6.10}{[1][35][]35}}
\newlabel{eje:6-unif-lgn}{{6.11}{35}{}{ejercicio.6.11}{}}
\newlabel{eje:6-unif-lgn@cref}{{[ejercicio][11][6]6.11}{[1][35][]35}}
\newlabel{eje:6-encuesta}{{6.12}{35}{}{ejercicio.6.12}{}}
\newlabel{eje:6-encuesta@cref}{{[ejercicio][12][6]6.12}{[1][35][]35}}
\newlabel{eje:6-cheb-exacta}{{6.13}{36}{}{ejercicio.6.13}{}}
\newlabel{eje:6-cheb-exacta@cref}{{[ejercicio][13][6]6.13}{[1][36][]36}}
\newlabel{eje:6-cantelli}{{6.14}{36}{}{ejercicio.6.14}{}}
\newlabel{eje:6-cantelli@cref}{{[ejercicio][14][6]6.14}{[1][36][]36}}
\newlabel{eje:6-expo}{{6.15}{36}{}{ejercicio.6.15}{}}
\newlabel{eje:6-expo@cref}{{[ejercicio][15][6]6.15}{[1][36][]36}}
\newlabel{eje:6-afin-normal}{{6.16}{36}{}{ejercicio.6.16}{}}
\newlabel{eje:6-afin-normal@cref}{{[ejercicio][16][6]6.16}{[1][36][]36}}
\newlabel{ap:soluciones}{{A}{37}{Soluciones de los ejercicios}{appendix.A}{}}
\newlabel{ap:soluciones@cref}{{[appendix][1][2147483647]A}{[1][37][]37}}
\newlabel{sol:eje:1-reglas}{{A.1}{37}{Tema 1}{section*.3}{}}
\newlabel{sol:eje:1-reglas@cref}{{[subappendix][1][2147483647,1]A.1}{[1][37][]37}}
\newlabel{sol:eje:1-vocabulario}{{A.1}{37}{Tema 1}{section*.4}{}}
\newlabel{sol:eje:1-vocabulario@cref}{{[subappendix][1][2147483647,1]A.1}{[1][37][]37}}
\newlabel{sol:eje:1-contrasenas}{{A.1}{38}{Tema 1}{section*.5}{}}
\newlabel{sol:eje:1-contrasenas@cref}{{[subappendix][1][2147483647,1]A.1}{[1][38][]38}}
\newlabel{sol:eje:1-divisibles}{{A.1}{38}{Tema 1}{section*.6}{}}
\newlabel{sol:eje:1-divisibles@cref}{{[subappendix][1][2147483647,1]A.1}{[1][38][]38}}
\newlabel{sol:eje:1-palindromos}{{A.1}{38}{Tema 1}{section*.7}{}}
\newlabel{sol:eje:1-palindromos@cref}{{[subappendix][1][2147483647,1]A.1}{[1][38][]38}}
\newlabel{sol:eje:1-mesa}{{A.1}{39}{Tema 1}{section*.8}{}}
\newlabel{sol:eje:1-mesa@cref}{{[subappendix][1][2147483647,1]A.1}{[1][39][]39}}
\newlabel{sol:eje:1-cifras}{{A.1}{39}{Tema 1}{section*.9}{}}
\newlabel{sol:eje:1-cifras@cref}{{[subappendix][1][2147483647,1]A.1}{[1][39][]39}}
\newlabel{sol:eje:1-funciones}{{A.1}{39}{Tema 1}{section*.10}{}}
\newlabel{sol:eje:1-funciones@cref}{{[subappendix][1][2147483647,1]A.1}{[1][39][]39}}
\newlabel{sol:eje:1-sombreros}{{A.1}{40}{Tema 1}{section*.11}{}}
\newlabel{sol:eje:1-sombreros@cref}{{[subappendix][1][2147483647,1]A.1}{[1][40][]40}}
\newlabel{sol:eje:1-encuesta}{{A.1}{40}{Tema 1}{section*.12}{}}
\newlabel{sol:eje:1-encuesta@cref}{{[subappendix][1][2147483647,1]A.1}{[1][40][]40}}
\newlabel{sol:eje:1-producto-ind}{{A.1}{41}{Tema 1}{section*.13}{}}
\newlabel{sol:eje:1-producto-ind@cref}{{[subappendix][1][2147483647,1]A.1}{[1][40][]41}}
\newlabel{sol:eje:1-fibonacci}{{A.1}{41}{Tema 1}{section*.14}{}}
\newlabel{sol:eje:1-fibonacci@cref}{{[subappendix][1][2147483647,1]A.1}{[1][41][]41}}
\newlabel{sol:eje:1-ie-cuatro}{{A.1}{41}{Tema 1}{section*.15}{}}
\newlabel{sol:eje:1-ie-cuatro@cref}{{[subappendix][1][2147483647,1]A.1}{[1][41][]41}}
\newlabel{sol:eje:1-collares}{{A.1}{42}{Tema 1}{section*.16}{}}
\newlabel{sol:eje:1-collares@cref}{{[subappendix][1][2147483647,1]A.1}{[1][42][]42}}
\newlabel{sol:eje:2-calcular}{{A.2}{42}{Tema 2}{section*.17}{}}
\newlabel{sol:eje:2-calcular@cref}{{[subappendix][2][2147483647,1]A.2}{[1][42][]42}}
\newlabel{sol:eje:2-modelo}{{A.2}{43}{Tema 2}{section*.18}{}}
\newlabel{sol:eje:2-modelo@cref}{{[subappendix][2][2147483647,1]A.2}{[1][43][]43}}
\newlabel{sol:eje:2-poker}{{A.2}{43}{Tema 2}{section*.19}{}}
\newlabel{sol:eje:2-poker@cref}{{[subappendix][2][2147483647,1]A.2}{[1][43][]43}}
\newlabel{sol:eje:2-equipos}{{A.2}{43}{Tema 2}{section*.20}{}}
\newlabel{sol:eje:2-equipos@cref}{{[subappendix][2][2147483647,1]A.2}{[1][43][]43}}
\newlabel{sol:eje:2-anagramas}{{A.2}{44}{Tema 2}{section*.21}{}}
\newlabel{sol:eje:2-anagramas@cref}{{[subappendix][2][2147483647,1]A.2}{[1][44][]44}}
\newlabel{sol:eje:2-coeficientes}{{A.2}{44}{Tema 2}{section*.22}{}}
\newlabel{sol:eje:2-coeficientes@cref}{{[subappendix][2][2147483647,1]A.2}{[1][44][]44}}
\newlabel{sol:eje:2-ecuaciones}{{A.2}{44}{Tema 2}{section*.23}{}}
\newlabel{sol:eje:2-ecuaciones@cref}{{[subappendix][2][2147483647,1]A.2}{[1][44][]44}}
\newlabel{sol:eje:2-caminos}{{A.2}{45}{Tema 2}{section*.24}{}}
\newlabel{sol:eje:2-caminos@cref}{{[subappendix][2][2147483647,1]A.2}{[1][45][]45}}
\newlabel{sol:eje:2-estaciones}{{A.2}{45}{Tema 2}{section*.25}{}}
\newlabel{sol:eje:2-estaciones@cref}{{[subappendix][2][2147483647,1]A.2}{[1][45][]45}}
\newlabel{sol:eje:2-desarreglos}{{A.2}{45}{Tema 2}{section*.26}{}}
\newlabel{sol:eje:2-desarreglos@cref}{{[subappendix][2][2147483647,1]A.2}{[1][45][]45}}
\newlabel{sol:eje:2-ositos}{{A.2}{46}{Tema 2}{section*.27}{}}
\newlabel{sol:eje:2-ositos@cref}{{[subappendix][2][2147483647,1]A.2}{[1][46][]46}}
\newlabel{sol:eje:2-palo}{{A.2}{46}{Tema 2}{section*.28}{}}
\newlabel{sol:eje:2-palo@cref}{{[subappendix][2][2147483647,1]A.2}{[1][46][]46}}
\newlabel{sol:eje:2-sumak}{{A.2}{46}{Tema 2}{section*.29}{}}
\newlabel{sol:eje:2-sumak@cref}{{[subappendix][2][2147483647,1]A.2}{[1][46][]46}}
\newlabel{sol:eje:2-vandermonde}{{A.2}{47}{Tema 2}{section*.30}{}}
\newlabel{sol:eje:2-vandermonde@cref}{{[subappendix][2][2147483647,1]A.2}{[1][47][]47}}
\newlabel{sol:eje:2-pares}{{A.2}{47}{Tema 2}{section*.31}{}}
\newlabel{sol:eje:2-pares@cref}{{[subappendix][2][2147483647,1]A.2}{[1][47][]47}}
\newlabel{sol:eje:3-espacio}{{A.3}{48}{Tema 3}{section*.32}{}}
\newlabel{sol:eje:3-espacio@cref}{{[subappendix][3][2147483647,1]A.3}{[1][47][]48}}
\newlabel{sol:eje:3-sucesos}{{A.3}{48}{Tema 3}{section*.33}{}}
\newlabel{sol:eje:3-sucesos@cref}{{[subappendix][3][2147483647,1]A.3}{[1][48][]48}}
\newlabel{sol:eje:3-monedas}{{A.3}{48}{Tema 3}{section*.34}{}}
\newlabel{sol:eje:3-monedas@cref}{{[subappendix][3][2147483647,1]A.3}{[1][48][]48}}
\newlabel{sol:eje:3-dados}{{A.3}{49}{Tema 3}{section*.35}{}}
\newlabel{sol:eje:3-dados@cref}{{[subappendix][3][2147483647,1]A.3}{[1][49][]49}}
\newlabel{sol:eje:3-cuatro-dados}{{A.3}{49}{Tema 3}{section*.36}{}}
\newlabel{sol:eje:3-cuatro-dados@cref}{{[subappendix][3][2147483647,1]A.3}{[1][49][]49}}
\newlabel{sol:eje:3-horarios}{{A.3}{49}{Tema 3}{section*.37}{}}
\newlabel{sol:eje:3-horarios@cref}{{[subappendix][3][2147483647,1]A.3}{[1][49][]49}}
\newlabel{sol:eje:3-robos}{{A.3}{50}{Tema 3}{section*.38}{}}
\newlabel{sol:eje:3-robos@cref}{{[subappendix][3][2147483647,1]A.3}{[1][49][]50}}
\newlabel{sol:eje:3-seis-cartas}{{A.3}{50}{Tema 3}{section*.39}{}}
\newlabel{sol:eje:3-seis-cartas@cref}{{[subappendix][3][2147483647,1]A.3}{[1][50][]50}}
\newlabel{sol:eje:3-cumple}{{A.3}{50}{Tema 3}{section*.40}{}}
\newlabel{sol:eje:3-cumple@cref}{{[subappendix][3][2147483647,1]A.3}{[1][50][]50}}
\newlabel{sol:eje:3-urna}{{A.3}{51}{Tema 3}{section*.41}{}}
\newlabel{sol:eje:3-urna@cref}{{[subappendix][3][2147483647,1]A.3}{[1][50][]51}}
\newlabel{sol:eje:3-geometrica}{{A.3}{51}{Tema 3}{section*.42}{}}
\newlabel{sol:eje:3-geometrica@cref}{{[subappendix][3][2147483647,1]A.3}{[1][51][]51}}
\newlabel{sol:eje:3-bonferroni}{{A.3}{51}{Tema 3}{section*.43}{}}
\newlabel{sol:eje:3-bonferroni@cref}{{[subappendix][3][2147483647,1]A.3}{[1][51][]51}}
\newlabel{sol:eje:3-interseccion}{{A.3}{51}{Tema 3}{section*.44}{}}
\newlabel{sol:eje:3-interseccion@cref}{{[subappendix][3][2147483647,1]A.3}{[1][51][]51}}
\newlabel{sol:eje:3-distancia}{{A.3}{52}{Tema 3}{section*.45}{}}
\newlabel{sol:eje:3-distancia@cref}{{[subappendix][3][2147483647,1]A.3}{[1][52][]52}}
\newlabel{sol:eje:3-hash}{{A.3}{52}{Tema 3}{section*.46}{}}
\newlabel{sol:eje:3-hash@cref}{{[subappendix][3][2147483647,1]A.3}{[1][52][]52}}
\newlabel{sol:eje:3-laplace}{{A.3}{52}{Tema 3}{section*.47}{}}
\newlabel{sol:eje:3-laplace@cref}{{[subappendix][3][2147483647,1]A.3}{[1][52][]52}}
\newlabel{sol:eje:4-tabla}{{A.4}{53}{Tema 4}{section*.48}{}}
\newlabel{sol:eje:4-tabla@cref}{{[subappendix][4][2147483647,1]A.4}{[1][53][]53}}
\newlabel{sol:eje:4-cachorros}{{A.4}{53}{Tema 4}{section*.49}{}}
\newlabel{sol:eje:4-cachorros@cref}{{[subappendix][4][2147483647,1]A.4}{[1][53][]53}}
\newlabel{sol:eje:4-spam}{{A.4}{54}{Tema 4}{section*.50}{}}
\newlabel{sol:eje:4-spam@cref}{{[subappendix][4][2147483647,1]A.4}{[1][53][]54}}
\newlabel{sol:eje:4-dos-pruebas}{{A.4}{54}{Tema 4}{section*.51}{}}
\newlabel{sol:eje:4-dos-pruebas@cref}{{[subappendix][4][2147483647,1]A.4}{[1][54][]54}}
\newlabel{sol:eje:4-piezas}{{A.4}{54}{Tema 4}{section*.52}{}}
\newlabel{sol:eje:4-piezas@cref}{{[subappendix][4][2147483647,1]A.4}{[1][54][]54}}
\newlabel{sol:eje:4-perro}{{A.4}{55}{Tema 4}{section*.53}{}}
\newlabel{sol:eje:4-perro@cref}{{[subappendix][4][2147483647,1]A.4}{[1][55][]55}}
\newlabel{sol:eje:4-monty7}{{A.4}{55}{Tema 4}{section*.54}{}}
\newlabel{sol:eje:4-monty7@cref}{{[subappendix][4][2147483647,1]A.4}{[1][55][]55}}
\newlabel{sol:eje:4-monty-p}{{A.4}{55}{Tema 4}{section*.55}{}}
\newlabel{sol:eje:4-monty-p@cref}{{[subappendix][4][2147483647,1]A.4}{[1][55][]55}}
\newlabel{sol:eje:4-paridad}{{A.4}{56}{Tema 4}{section*.56}{}}
\newlabel{sol:eje:4-paridad@cref}{{[subappendix][4][2147483647,1]A.4}{[1][56][]56}}
\newlabel{sol:eje:4-ajedrez}{{A.4}{56}{Tema 4}{section*.57}{}}
\newlabel{sol:eje:4-ajedrez@cref}{{[subappendix][4][2147483647,1]A.4}{[1][56][]56}}
\newlabel{sol:eje:4-si-mismo}{{A.4}{56}{Tema 4}{section*.58}{}}
\newlabel{sol:eje:4-si-mismo@cref}{{[subappendix][4][2147483647,1]A.4}{[1][56][]56}}
\newlabel{sol:eje:4-union-indep}{{A.4}{57}{Tema 4}{section*.59}{}}
\newlabel{sol:eje:4-union-indep@cref}{{[subappendix][4][2147483647,1]A.4}{[1][57][]57}}
\newlabel{sol:eje:4-producto-triple}{{A.4}{57}{Tema 4}{section*.60}{}}
\newlabel{sol:eje:4-producto-triple@cref}{{[subappendix][4][2147483647,1]A.4}{[1][57][]57}}
\newlabel{sol:eje:4-razones}{{A.4}{58}{Tema 4}{section*.61}{}}
\newlabel{sol:eje:4-razones@cref}{{[subappendix][4][2147483647,1]A.4}{[1][57][]58}}
\newlabel{sol:eje:4-simetria}{{A.4}{58}{Tema 4}{section*.62}{}}
\newlabel{sol:eje:4-simetria@cref}{{[subappendix][4][2147483647,1]A.4}{[1][58][]58}}
\newlabel{sol:eje:5-masa}{{A.5}{58}{Tema 5}{section*.63}{}}
\newlabel{sol:eje:5-masa@cref}{{[subappendix][5][2147483647,1]A.5}{[1][58][]58}}
\newlabel{sol:eje:5-diferencia}{{A.5}{59}{Tema 5}{section*.64}{}}
\newlabel{sol:eje:5-diferencia@cref}{{[subappendix][5][2147483647,1]A.5}{[1][59][]59}}
\newlabel{sol:eje:5-binomial}{{A.5}{59}{Tema 5}{section*.65}{}}
\newlabel{sol:eje:5-binomial@cref}{{[subappendix][5][2147483647,1]A.5}{[1][59][]59}}
\newlabel{sol:eje:5-geometrica}{{A.5}{59}{Tema 5}{section*.66}{}}
\newlabel{sol:eje:5-geometrica@cref}{{[subappendix][5][2147483647,1]A.5}{[1][59][]59}}
\newlabel{sol:eje:5-poisson-aprox}{{A.5}{60}{Tema 5}{section*.67}{}}
\newlabel{sol:eje:5-poisson-aprox@cref}{{[subappendix][5][2147483647,1]A.5}{[1][60][]60}}
\newlabel{sol:eje:5-cumple-esp}{{A.5}{60}{Tema 5}{section*.68}{}}
\newlabel{sol:eje:5-cumple-esp@cref}{{[subappendix][5][2147483647,1]A.5}{[1][60][]60}}
\newlabel{sol:eje:5-hash}{{A.5}{60}{Tema 5}{section*.69}{}}
\newlabel{sol:eje:5-hash@cref}{{[subappendix][5][2147483647,1]A.5}{[1][60][]60}}
\newlabel{sol:eje:5-hijos}{{A.5}{61}{Tema 5}{section*.70}{}}
\newlabel{sol:eje:5-hijos@cref}{{[subappendix][5][2147483647,1]A.5}{[1][61][]61}}
\newlabel{sol:eje:5-inversiones}{{A.5}{61}{Tema 5}{section*.71}{}}
\newlabel{sol:eje:5-inversiones@cref}{{[subappendix][5][2147483647,1]A.5}{[1][61][]61}}
\newlabel{sol:eje:5-categorica}{{A.5}{61}{Tema 5}{section*.72}{}}
\newlabel{sol:eje:5-categorica@cref}{{[subappendix][5][2147483647,1]A.5}{[1][61][]61}}
\newlabel{sol:eje:5-passk}{{A.5}{61}{Tema 5}{section*.73}{}}
\newlabel{sol:eje:5-passk@cref}{{[subappendix][5][2147483647,1]A.5}{[1][61][]61}}
\newlabel{sol:eje:5-colas}{{A.5}{62}{Tema 5}{section*.74}{}}
\newlabel{sol:eje:5-colas@cref}{{[subappendix][5][2147483647,1]A.5}{[1][62][]62}}
\newlabel{sol:eje:5-poisson}{{A.5}{62}{Tema 5}{section*.75}{}}
\newlabel{sol:eje:5-poisson@cref}{{[subappendix][5][2147483647,1]A.5}{[1][62][]62}}
\newlabel{sol:eje:5-sin-memoria}{{A.5}{62}{Tema 5}{section*.76}{}}
\newlabel{sol:eje:5-sin-memoria@cref}{{[subappendix][5][2147483647,1]A.5}{[1][62][]62}}
\newlabel{sol:eje:5-geom-cond}{{A.5}{63}{Tema 5}{section*.77}{}}
\newlabel{sol:eje:5-geom-cond@cref}{{[subappendix][5][2147483647,1]A.5}{[1][63][]63}}
\newlabel{sol:eje:6-conceptos}{{A.6}{63}{Tema 6}{section*.78}{}}
\newlabel{sol:eje:6-conceptos@cref}{{[subappendix][6][2147483647,1]A.6}{[1][63][]63}}
\newlabel{sol:eje:6-densidad}{{A.6}{64}{Tema 6}{section*.79}{}}
\newlabel{sol:eje:6-densidad@cref}{{[subappendix][6][2147483647,1]A.6}{[1][63][]64}}
\newlabel{sol:eje:6-pasas}{{A.6}{64}{Tema 6}{section*.80}{}}
\newlabel{sol:eje:6-pasas@cref}{{[subappendix][6][2147483647,1]A.6}{[1][64][]64}}
\newlabel{sol:eje:6-exponencial}{{A.6}{64}{Tema 6}{section*.81}{}}
\newlabel{sol:eje:6-exponencial@cref}{{[subappendix][6][2147483647,1]A.6}{[1][64][]64}}
\newlabel{sol:eje:6-autobus}{{A.6}{64}{Tema 6}{section*.82}{}}
\newlabel{sol:eje:6-autobus@cref}{{[subappendix][6][2147483647,1]A.6}{[1][64][]64}}
\newlabel{sol:eje:6-tipificar}{{A.6}{65}{Tema 6}{section*.83}{}}
\newlabel{sol:eje:6-tipificar@cref}{{[subappendix][6][2147483647,1]A.6}{[1][65][]65}}
\newlabel{sol:eje:6-gestacion}{{A.6}{65}{Tema 6}{section*.84}{}}
\newlabel{sol:eje:6-gestacion@cref}{{[subappendix][6][2147483647,1]A.6}{[1][65][]65}}
\newlabel{sol:eje:6-palo}{{A.6}{65}{Tema 6}{section*.85}{}}
\newlabel{sol:eje:6-palo@cref}{{[subappendix][6][2147483647,1]A.6}{[1][65][]65}}
\newlabel{sol:eje:6-markov}{{A.6}{66}{Tema 6}{section*.86}{}}
\newlabel{sol:eje:6-markov@cref}{{[subappendix][6][2147483647,1]A.6}{[1][65][]66}}
\newlabel{sol:eje:6-n-cheb}{{A.6}{66}{Tema 6}{section*.87}{}}
\newlabel{sol:eje:6-n-cheb@cref}{{[subappendix][6][2147483647,1]A.6}{[1][66][]66}}
\newlabel{sol:eje:6-unif-lgn}{{A.6}{66}{Tema 6}{section*.88}{}}
\newlabel{sol:eje:6-unif-lgn@cref}{{[subappendix][6][2147483647,1]A.6}{[1][66][]66}}
\newlabel{sol:eje:6-encuesta}{{A.6}{66}{Tema 6}{section*.89}{}}
\newlabel{sol:eje:6-encuesta@cref}{{[subappendix][6][2147483647,1]A.6}{[1][66][]66}}
\newlabel{sol:eje:6-cheb-exacta}{{A.6}{67}{Tema 6}{section*.90}{}}
\newlabel{sol:eje:6-cheb-exacta@cref}{{[subappendix][6][2147483647,1]A.6}{[1][67][]67}}
\newlabel{sol:eje:6-cantelli}{{A.6}{67}{Tema 6}{section*.91}{}}
\newlabel{sol:eje:6-cantelli@cref}{{[subappendix][6][2147483647,1]A.6}{[1][67][]67}}
\newlabel{sol:eje:6-expo}{{A.6}{67}{Tema 6}{section*.92}{}}
\newlabel{sol:eje:6-expo@cref}{{[subappendix][6][2147483647,1]A.6}{[1][67][]67}}
\newlabel{sol:eje:6-afin-normal}{{A.6}{68}{Tema 6}{section*.93}{}}
\newlabel{sol:eje:6-afin-normal@cref}{{[subappendix][6][2147483647,1]A.6}{[1][68][]68}}
\makeatother

\begin{document}
\frontmatter
\begin{titlepage}
\centering
\vspace*{2.5cm}
{\large\scshape Anatomía de un LLM}\par
\vspace{0.4cm}
{\normalsize Fundamentos matemáticos de los modelos de lenguaje}\par
\vspace{2.2cm}
\rule{\textwidth}{0.8pt}\par\vspace{0.5cm}
{\small\scshape Fase 0 \textperiodcentered\ Preliminares matemáticos}\par\vspace{0.3cm}
{\Huge\bfseries Módulo 0.8\par}\vspace{0.3cm}
{\LARGE Combinatoria y probabilidad\par}
\vspace{0.5cm}\rule{\textwidth}{0.4pt}\par
\vspace{1.6cm}
{\large Resumen del módulo}\par
\vfill
{\small Edición de octubre de 2026}\par
\end{titlepage}

\chapter*{Sobre este resumen}
\addcontentsline{toc}{chapter}{Sobre este resumen}
\markboth{Sobre este resumen}{}

Este documento reúne lo indispensable de los apuntes del Módulo 0.8: el texto expositivo, todas las definiciones y todos los resultados (teoremas, proposiciones, lemas y corolarios), las tablas y las observaciones de contenido. Se han suprimido las demostraciones, los ejemplos, las lecturas, las figuras, los ejercicios y sus soluciones. La numeración de temas, secciones, enunciados y ecuaciones es la misma que en los apuntes, de modo que cualquier referencia puede buscarse allí; las que remiten a elementos suprimidos (un ejemplo, una figura, una ecuación de una demostración) conservan su número en los apuntes.

\section*{Notación}
\begin{center}
\begin{tabular}{@{}ll@{}}
\toprule
Símbolo & Significado \\
\midrule
$\abs{A}$, $[n]$ & cardinal de un conjunto finito; $[n]=\{1,\dots,n\}$ \\
$n!$, $n^{\underline{k}}$ & factorial; factorial descendente $n(n-1)\cdots(n-k+1)$ \\
$\binom nk$, $\binom{n}{k_1,\dots,k_r}$ & número combinatorio; coeficiente multinomial \\
$\Omega$, $A^{c}$, $A\dsim B$ & espacio muestral; suceso contrario; diferencia simétrica \\
$\PP(A)$, $\PP(A\mid B)$ & probabilidad; probabilidad condicionada \\
$I_A$ & variable indicadora del suceso $A$ \\
$p_X$, $f_X$, $F_X$ & función de masa, densidad y función de distribución de $X$ \\
$\E[X]$, $\Var(X)$, $\sigma_X$ & esperanza, varianza, desviación típica \\
$\Ber(p)$, $\Bin(n,p)$, $\Geom(p)$ & distribuciones discretas \\
$\Poi(\lambda)$, $\Cat(\bm p)$ & distribuciones de Poisson y categórica \\
$\Unif(a,b)$, $\Expo(\lambda)$, $\Nor(\mu,\sigma^{2})$ & distribuciones continuas \\
$\varphi$, $\Phi$ & densidad y función de distribución de la normal estándar \\
$\Xbar_n$ & media muestral $\frac1n\sum_{i=1}^{n}X_i$ \\
\bottomrule
\end{tabular}
\end{center}


\tableofcontents

\mainmatter
% =====================================================================
\chapter{Principios de conteo y ordenaciones}
\label{tema:1}
% =====================================================================

Contar bien es el primer paso para calcular probabilidades en espacios finitos. Este tema reúne las reglas básicas del recuento (suma, producto, complemento, biyección y división), el principio de inclusión--exclusión para dos y tres conjuntos y el recuento de ordenaciones: variaciones y permutaciones, con y sin repetición. La presentación sigue el capítulo~14, «Cardinality Rules», de Lehman, Leighton y Meyer~\cite{mcs}, texto del curso 6.042J del MIT; el capítulo~1 de Ross~\cite{ross} (§1.2 y §1.3); la sección~1.4, «How to count», de Blitzstein y Hwang~\cite{bh}, y la sección~3.1 de Grinstead y Snell~\cite{gs}.

\section{Reglas de la suma, del producto y del complemento}

Se recuerda que, para un conjunto finito $A$, $\abs{A}$ denota su número de elementos y que $[n]=\{1,2,\dots,n\}$. Contar un conjunto consiste en hallar $\abs{A}$; casi siempre se hace descomponiendo $A$ en piezas más sencillas o poniéndolo en correspondencia con otro conjunto ya contado.

\begin{proposicion}[Regla de la suma]\label{prop:t1-suma}
Si $A_1,\dots,A_m$ son conjuntos finitos disjuntos dos a dos, entonces
\[
\abs{A_1\cup\dots\cup A_m}=\abs{A_1}+\dots+\abs{A_m}.
\]
\end{proposicion}

\begin{proposicion}[Regla del producto]\label{prop:t1-producto}
Si $A_1,\dots,A_k$ son conjuntos finitos, $\abs{A_1\times\dots\times A_k}=\abs{A_1}\cdots\abs{A_k}$.
\end{proposicion}

En la práctica se usa una versión más flexible, en la que el conjunto de opciones de cada etapa puede depender de las elecciones previas, pero no su número.

\begin{teorema}[Regla del producto generalizada]\label{teo:t1-producto-gen}
Sea $S$ un conjunto de sucesiones $(s_1,\dots,s_k)$ tal que hay $n_1$ valores posibles de $s_1$ y, para cada elección de $s_1,\dots,s_{i-1}$, hay exactamente $n_i$ valores posibles de $s_i$ ($i=2,\dots,k$). Entonces $\abs{S}=n_1n_2\cdots n_k$.
\end{teorema}

\begin{proposicion}[Regla del complemento]\label{prop:t1-complemento}
Si $A\subseteq U$ y $U$ es finito, $\abs{U\setminus A}=\abs{U}-\abs{A}$.
\end{proposicion}

\addtocounter{definicion}{1}%

\section{Reglas de la biyección y de la división}

\begin{proposicion}[Regla de la biyección]\label{prop:t1-biyeccion}
Si existe una biyección $f\colon A\to B$ entre conjuntos finitos, entonces $\abs{A}=\abs{B}$.
\end{proposicion}

La demostración es inmediata: si $A=\{a_1,\dots,a_n\}$ sin repeticiones, entonces $f(a_1),\dots,f(a_n)$ son distintos (inyectividad) y agotan $B$ (sobreyectividad). Esta regla permite contar un conjunto contando otro~\cite[§14.1]{mcs}.

\addtocounter{definicion}{1}%

\begin{proposicion}[Regla de la división]\label{prop:t1-division}
Sea $f\colon A\to B$ una función sobreyectiva entre conjuntos finitos tal que cada $b\in B$ tiene exactamente $d$ preimágenes ($f$ es \emph{$d$ a~1}). Entonces $\abs{A}=d\,\abs{B}$.
\end{proposicion}

\section{Inclusión--exclusión para dos y tres conjuntos}

Cuando los conjuntos no son disjuntos, la suma de sus cardinales cuenta dos veces los elementos comunes y hay que corregirla.

\begin{teorema}[Inclusión--exclusión]\label{teo:t1-ie}
Para conjuntos finitos $A$, $B$ y $C$:
\begin{align*}
\abs{A\cup B}&=\abs{A}+\abs{B}-\abs{A\cap B},\\
\abs{A\cup B\cup C}&=\abs{A}+\abs{B}+\abs{C}-\abs{A\cap B}-\abs{A\cap C}-\abs{B\cap C}+\abs{A\cap B\cap C}.
\end{align*}
\end{teorema}

La versión para $m$ conjuntos se demuestra en el \cref{tema:2}, una vez disponibles los números combinatorios.

\addtocounter{definicion}{1}%

\section{Variaciones y permutaciones}

Se trata ahora de contar \emph{selecciones ordenadas} de elementos de un conjunto de $n$ elementos. La terminología española distingue según se permita o no repetir.

\begin{definicion}[Variaciones y permutaciones]
Sea $A$ un conjunto con $n$ elementos y $k\ge 0$ un entero.
\begin{enumerate}
  \item Una \emph{variación con repetición} de orden $k$ es una sucesión $(a_1,\dots,a_k)$ de elementos de $A$, posiblemente repetidos; equivale a una función $[k]\to A$.
  \item Una \emph{variación sin repetición} de orden $k\le n$ es una sucesión de $k$ elementos \emph{distintos} de $A$; equivale a una función inyectiva $[k]\to A$.
  \item Una \emph{permutación} de $A$ es una variación sin repetición de orden $n$, es decir, una ordenación de todos sus elementos; equivale a una biyección $[n]\to A$.
\end{enumerate}
Se escribe $n!=1\cdot2\cdots n$, con $0!=1$, y $n^{\underline{k}}=n(n-1)\cdots(n-k+1)$ (\emph{factorial descendente}, con $n^{\underline{0}}=1$).
\end{definicion}

\begin{teorema}\label{teo:t1-variaciones}
Si $\abs{A}=n$, el número de variaciones con repetición de orden $k$ es $n^{k}$; el de variaciones sin repetición de orden $k\le n$ es $n^{\underline{k}}=n!/(n-k)!$, y el de permutaciones es $n!$.
\end{teorema}

\addtocounter{definicion}{1}%

\addtocounter{definicion}{1}%

\addtocounter{definicion}{1}%

\addtocounter{definicion}{1}%


% =====================================================================
\chapter{Combinaciones, números combinatorios y doble conteo}
\label{tema:2}
% =====================================================================

Este tema cuenta \emph{selecciones no ordenadas}: subconjuntos, multiconjuntos y repartos. Aparecen los números combinatorios, el binomio y el multinomio de Newton, la técnica del doble conteo y la forma general del principio de inclusión--exclusión. Se siguen las secciones §14.5 a §14.7 de Lehman, Leighton y Meyer~\cite{mcs}; las secciones §1.4 a §1.6 de Ross~\cite{ross}; las secciones §1.4 y §1.5, «Story proofs», de Blitzstein y Hwang~\cite{bh}, cuyo estilo de «demostración mediante una historia» se adopta aquí, y la sección~3.2 de Grinstead y Snell~\cite{gs}.

\section{Combinaciones y números combinatorios}

\begin{definicion}[Combinaciones]
Sea $A$ un conjunto con $n$ elementos y $0\le k\le n$. Una \emph{combinación} (sin repetición) de orden $k$ de $A$ es un subconjunto de $A$ con $k$ elementos. Su número se denota $\binom{n}{k}$ y se llama \emph{número combinatorio} («$n$ sobre $k$»). Por convenio, $\binom{n}{k}=0$ si $k<0$ o $k>n$.
\end{definicion}

\begin{teorema}\label{teo:t2-binom}
Para $0\le k\le n$,
\[
\binom{n}{k}=\frac{n^{\underline{k}}}{k!}=\frac{n!}{k!\,(n-k)!}.
\]
\end{teorema}

\begin{proposicion}[Propiedades básicas]\label{prop:t2-props}
Para enteros $n\ge1$ y $k$:
\begin{enumerate}
  \item simetría: $\binom{n}{k}=\binom{n}{n-k}$;
  \item fórmula de Pascal: $\binom{n}{k}=\binom{n-1}{k}+\binom{n-1}{k-1}$;
  \item absorción: $k\binom{n}{k}=n\binom{n-1}{k-1}$.
\end{enumerate}
\end{proposicion}

La fórmula de Pascal permite construir el \emph{triángulo de Pascal}, en el que cada número es la suma de los dos que tiene encima:
\[
\begin{array}{c}
1\\
1\quad 1\\
1\quad 2\quad 1\\
1\quad 3\quad 3\quad 1\\
1\quad 4\quad 6\quad 4\quad 1\\
1\quad 5\quad 10\quad 10\quad 5\quad 1
\end{array}
\]

\section{Binomio y multinomio de Newton}

\begin{teorema}[Binomio de Newton]\label{teo:t2-binomio}
Para $x,y$ reales y $n\ge0$ entero,
\[
(x+y)^{n}=\sum_{k=0}^{n}\binom{n}{k}x^{k}y^{n-k}.
\]
\end{teorema}

El módulo~0.3 demuestra este teorema por inducción; la demostración anterior es la combinatoria, que explica \emph{por qué} aparecen los números combinatorios.

\begin{corolario}\label{cor:t2-sumas}
Para $n\ge 0$, $\sum_{k=0}^{n}\binom{n}{k}=2^{n}$, y para $n\ge1$, $\sum_{k=0}^{n}(-1)^{k}\binom{n}{k}=0$.
\end{corolario}

\begin{definicion}[Coeficiente multinomial]
Sean $n=k_1+\dots+k_r$ con $k_i\ge0$ enteros. Se define
\[
\binom{n}{k_1,k_2,\dots,k_r}=\frac{n!}{k_1!\,k_2!\cdots k_r!}.
\]
\end{definicion}

\begin{teorema}[Permutaciones con repetición]\label{teo:t2-multinomial}
El número de sucesiones de longitud $n$ formadas con $k_1$ copias del símbolo $s_1$, $k_2$ del símbolo $s_2$, \dots, $k_r$ del símbolo $s_r$ es $\binom{n}{k_1,\dots,k_r}$. Equivalentemente, es el número de formas de repartir $n$ objetos distintos en $r$ cajas distinguibles, con $k_i$ objetos en la caja $i$.
\end{teorema}

\addtocounter{definicion}{1}%

\begin{teorema}[Multinomio de Newton]\label{teo:t2-multinomio}
Para reales $x_1,\dots,x_r$ y $n\ge0$,
\[
(x_1+\dots+x_r)^{n}=\sum_{\substack{k_1+\dots+k_r=n\\ k_i\ge0}}\binom{n}{k_1,\dots,k_r}x_1^{k_1}\cdots x_r^{k_r}.
\]
\end{teorema}

\section{Combinaciones con repetición}

\begin{definicion}
Una \emph{combinación con repetición} de orden $k$ de un conjunto de $n$ elementos es una selección no ordenada de $k$ elementos en la que se permite repetir; equivale a un vector $(x_1,\dots,x_n)$ de enteros $x_i\ge 0$ con $x_1+\dots+x_n=k$, donde $x_i$ es el número de veces que se elige el elemento $i$.
\end{definicion}

\begin{teorema}[Estrellas y barras]\label{teo:t2-estrellas}
El número de soluciones enteras de $x_1+\dots+x_n=k$ con todos los $x_i\ge0$ es $\binom{n+k-1}{k}=\binom{n+k-1}{n-1}$. El número de soluciones con todos los $x_i\ge1$ es $\binom{k-1}{n-1}$ (para $k\ge n\ge 1$).
\end{teorema}

\begin{table}[h]
\centering
\begin{tabular}{@{}lcc@{}}
\toprule
Selección de $k$ entre $n$ & sin repetición & con repetición \\
\midrule
ordenada (importa el orden) & $n^{\underline{k}}=\dfrac{n!}{(n-k)!}$ & $n^{k}$ \\[8pt]
no ordenada & $\dbinom{n}{k}$ & $\dbinom{n+k-1}{k}$ \\
\bottomrule
\end{tabular}
\caption{Los cuatro modelos básicos de selección. El caso con repetición y sin orden es el único en que la fórmula no es «evidente».}
\label{tab:t2-modelos}
\end{table}

\section{Doble conteo y demostraciones biyectivas}

Una \emph{demostración por doble conteo} establece una identidad $a=b$ mostrando que ambos lados cuentan el mismo conjunto de dos maneras distintas; una \emph{demostración biyectiva} la establece construyendo una biyección entre un conjunto de $a$ elementos y otro de $b$ elementos. Ya se han usado ambas en la \cref{prop:t2-props}. El esquema conviene escribirlo siempre de forma explícita: (1) describir el conjunto que se cuenta (una «historia»); (2) contarlo de una forma, obteniendo el lado izquierdo; (3) contarlo de otra, obteniendo el lado derecho~\cite[§1.5]{bh}.

\addtocounter{definicion}{1}%

\begin{proposicion}[Identidad de Vandermonde]\label{prop:t2-vandermonde}
Para enteros $m,n,k\ge0$, $\displaystyle\binom{m+n}{k}=\sum_{j=0}^{k}\binom{m}{j}\binom{n}{k-j}$.
\end{proposicion}

La demostración se pide en el \cref{eje:2-vandermonde}: un grupo de $m$ mujeres y $n$ hombres, y comités de $k$ personas clasificados por el número de mujeres.

\section{Inclusión--exclusión general}

\begin{teorema}[Inclusión--exclusión]\label{teo:t2-ie}
Para conjuntos finitos $A_1,\dots,A_m$,
\[
\Bigl|\bigcup_{i=1}^{m}A_i\Bigr|=\sum_{j=1}^{m}(-1)^{j+1}\sum_{1\le i_1<\dots<i_j\le m}\abs{A_{i_1}\cap\dots\cap A_{i_j}}.
\]
\end{teorema}

\addtocounter{definicion}{1}%

\addtocounter{definicion}{1}%

\addtocounter{definicion}{1}%


% =====================================================================
\chapter{Espacios de probabilidad}
\label{tema:3}
% =====================================================================

Este tema da la definición axiomática de probabilidad, deduce de los axiomas las reglas de cálculo y estudia el caso más importante para empezar, el modelo de Laplace, en el que calcular una probabilidad se reduce a contar. Se siguen la sección §1.2 de Bertsekas y Tsitsiklis~\cite{bt} (clase~1 de MIT 6.041~\cite{mit6041f10}), el capítulo~2 de Ross~\cite{ross} (§2.2 a §2.5), las secciones §1.2, §1.3 y §1.6 de Blitzstein y Hwang~\cite{bh} y el capítulo~16, «Events and Probability Spaces», de~\cite{mcs}, que incluye el problema del cumpleaños (§16.4).

\section{Experimentos, espacio muestral y sucesos}

\begin{definicion}[Espacio muestral y sucesos]
Un \emph{experimento aleatorio} es un procedimiento cuyo resultado no puede predecirse con certeza. El \emph{espacio muestral} $\Omega$ es el conjunto de todos sus resultados posibles. Un \emph{suceso} es un subconjunto $A\subseteq\Omega$; se dice que $A$ \emph{ocurre} si el resultado obtenido pertenece a $A$. Los sucesos $\{\omega\}$ se llaman \emph{elementales}.
\end{definicion}

Las operaciones con conjuntos del módulo~0.2 se traducen en operaciones con sucesos: $A\cup B$ es «ocurre $A$ o $B$»; $A\cap B$, «ocurren $A$ y $B$»; $A^{c}=\Omega\setminus A$, «no ocurre $A$»; $A\subseteq B$, «si ocurre $A$, ocurre $B$». Dos sucesos son \emph{incompatibles} (o mutuamente excluyentes) si $A\cap B=\emptyset$. Las leyes de De Morgan dan, por ejemplo, $(A\cup B)^{c}=A^{c}\cap B^{c}$: «no ocurre ninguno de los dos».

\addtocounter{definicion}{1}%

\section{Axiomas de Kolmogorov}

\begin{definicion}[Espacio de probabilidad]\label{def:t3-axiomas}
Una \emph{ley de probabilidad} sobre $\Omega$ es una función $\PP$ que asigna a cada suceso $A$ un número real $\PP(A)$ y que cumple:
\begin{enumerate}
  \item \emph{no negatividad}: $\PP(A)\ge0$ para todo suceso $A$;
  \item \emph{normalización}: $\PP(\Omega)=1$;
  \item \emph{aditividad numerable}: si $A_1,A_2,\dots$ son sucesos incompatibles dos a dos, $\PP\bigl(\bigcup_{i\ge1}A_i\bigr)=\sum_{i\ge1}\PP(A_i)$.
\end{enumerate}
El par $(\Omega,\PP)$ se llama \emph{espacio de probabilidad}.
\end{definicion}

\begin{observacion}
Si $\Omega$ es finito o numerable, todo subconjunto puede ser un suceso. Si $\Omega$ es no numerable, como un intervalo de $\R$, no es posible en general asignar probabilidad a todos los subconjuntos de forma coherente, y los sucesos se restringen a una familia adecuada (una $\sigma$-álgebra). En este módulo esa sutileza no interviene: todos los sucesos que aparecen (intervalos, uniones de intervalos, regiones del plano) son admisibles~\cite[§1.6]{bh}.
\end{observacion}

\section{Propiedades que se deducen de los axiomas}

\begin{proposicion}\label{prop:t3-props}
Sean $A$, $B$ y $A_1,\dots,A_n$ sucesos. Entonces:
\begin{enumerate}
  \item $\PP(\emptyset)=0$;
  \item \emph{aditividad finita}: si $A_1,\dots,A_n$ son incompatibles dos a dos, $\PP\bigl(\bigcup_{i=1}^{n}A_i\bigr)=\sum_{i=1}^{n}\PP(A_i)$;
  \item $\PP(A^{c})=1-\PP(A)$;
  \item si $A\subseteq B$, entonces $\PP(B\setminus A)=\PP(B)-\PP(A)$ y $\PP(A)\le\PP(B)$;
  \item $0\le\PP(A)\le1$;
  \item $\PP(A\cup B)=\PP(A)+\PP(B)-\PP(A\cap B)$.
\end{enumerate}
\end{proposicion}

\begin{teorema}[Desigualdad de Boole o cota de la unión]\label{teo:t3-boole}
Para sucesos cualesquiera $A_1,\dots,A_n$,
\(
\PP\bigl(\bigcup_{i=1}^{n}A_i\bigr)\le\sum_{i=1}^{n}\PP(A_i).
\)
\end{teorema}

La fórmula de inclusión--exclusión del \cref{teo:t2-ie} tiene una versión probabilística con la misma estructura; para tres sucesos,
\[
\PP(A\cup B\cup C)=\PP(A)+\PP(B)+\PP(C)-\PP(A\cap B)-\PP(A\cap C)-\PP(B\cap C)+\PP(A\cap B\cap C),
\]
y se demuestra a partir de (vi) exactamente como el \cref{teo:t1-ie}.

\section{Espacios discretos y modelo de Laplace}

\begin{proposicion}[Espacios discretos]\label{prop:t3-discreto}
Sea $\Omega$ finito o numerable y $p\colon\Omega\to[0,1]$ con $\sum_{\omega\in\Omega}p(\omega)=1$. Entonces $\PP(A)=\sum_{\omega\in A}p(\omega)$ define una ley de probabilidad. Recíprocamente, toda ley de probabilidad sobre $\Omega$ es de esta forma, con $p(\omega)=\PP(\{\omega\})$.
\end{proposicion}

\begin{definicion}[Modelo de Laplace]\label{def:t3-laplace}
Si $\Omega$ es finito y todos los resultados son equiprobables, $p(\omega)=1/\abs{\Omega}$, y entonces
\[
\PP(A)=\frac{\abs{A}}{\abs{\Omega}}=\frac{\text{casos favorables}}{\text{casos posibles}}.
\]
\end{definicion}

El modelo de Laplace es una \emph{hipótesis} sobre el experimento, no una consecuencia de los axiomas: es razonable para dados, monedas y cartas bien mezcladas, pero no, por ejemplo, para la suma de dos dados, cuyos 11 valores no son equiprobables. La técnica habitual consiste en elegir un espacio muestral en el que la equiprobabilidad sea evidente, normalmente uno de resultados \emph{ordenados} y distinguibles, aunque los objetos del problema parezcan indistinguibles.

\addtocounter{definicion}{1}%

\section{El problema del cumpleaños}

\begin{proposicion}\label{prop:t3-cumple}
Se asigna a cada una de $n$ personas un día entre $N$, de forma uniforme e independiente (todas las $N^{n}$ asignaciones son equiprobables). La probabilidad de que dos o más personas compartan día es
\[
p_n=1-\frac{N^{\underline{n}}}{N^{n}}=1-\prod_{i=0}^{n-1}\Bigl(1-\frac{i}{N}\Bigr),
\]
y además $1-p_n\le\exp\!\bigl(-\tfrac{n(n-1)}{2N}\bigr)$.
\end{proposicion}

\addtocounter{definicion}{1}%

\addtocounter{figure}{1}%

\addtocounter{definicion}{1}%

\addtocounter{definicion}{1}%


% =====================================================================
\chapter{Probabilidad condicionada, independencia y teorema de Bayes}
\label{tema:4}
% =====================================================================

La probabilidad condicionada formaliza cómo cambia la incertidumbre al recibir información. De ella se deducen la regla de la cadena, el teorema de la probabilidad total y el teorema de Bayes, y con ella se define la independencia. El tema sigue las secciones §1.3 a §1.5 de Bertsekas y Tsitsiklis~\cite{bt} (clases~2 y~3 de MIT 6.041~\cite{mit6041f10}), el capítulo~2 de Blitzstein y Hwang~\cite{bh}, el capítulo~3 de Ross~\cite{ross} y el capítulo~17 de~\cite{mcs}, cuyo tratamiento del problema de Monty Hall y de las pruebas médicas (§17.1 y §17.4.2) se adapta aquí.

\section{Probabilidad condicionada}

\begin{definicion}[Probabilidad condicionada]
Sean $A$ y $B$ sucesos con $\PP(B)>0$. La \emph{probabilidad de $A$ condicionada a $B$} es
\[
\PP(A\mid B)=\frac{\PP(A\cap B)}{\PP(B)}.
\]
\end{definicion}

La idea es que, al saber que ha ocurrido $B$, los resultados fuera de $B$ quedan descartados y los de $B$ conservan sus probabilidades \emph{relativas}; dividir por $\PP(B)$ restablece la normalización. En el modelo de Laplace, $\PP(A\mid B)=\abs{A\cap B}/\abs{B}$: se cuenta dentro de $B$ como si $B$ fuera el nuevo espacio muestral.

\begin{proposicion}\label{prop:t4-esprob}
Si $\PP(B)>0$, la función $A\mapsto\PP(A\mid B)$ es una ley de probabilidad sobre $\Omega$. En consecuencia, cumple todas las propiedades de la \cref{prop:t3-props}; por ejemplo, $\PP(A^{c}\mid B)=1-\PP(A\mid B)$.
\end{proposicion}

\addtocounter{definicion}{1}%

\section{Regla de la cadena}

\begin{teorema}[Regla de la cadena o de la multiplicación]\label{teo:t4-cadena}
Si $A_1,\dots,A_n$ son sucesos con $\PP(A_1\cap\dots\cap A_{n-1})>0$, entonces
\[
\PP(A_1\cap A_2\cap\dots\cap A_n)=\PP(A_1)\,\PP(A_2\mid A_1)\,\PP(A_3\mid A_1\cap A_2)\cdots\PP(A_n\mid A_1\cap\dots\cap A_{n-1}).
\]
\end{teorema}

\addtocounter{definicion}{1}%

\section{Probabilidad total y teorema de Bayes}

\begin{definicion}
Una familia de sucesos $B_1,\dots,B_n$ es una \emph{partición} de $\Omega$ si son incompatibles dos a dos y $B_1\cup\dots\cup B_n=\Omega$.
\end{definicion}

\begin{teorema}[Probabilidad total]\label{teo:t4-total}
Si $B_1,\dots,B_n$ es una partición de $\Omega$ con $\PP(B_i)>0$ para todo $i$, entonces, para todo suceso $A$,
\[
\PP(A)=\sum_{i=1}^{n}\PP(A\mid B_i)\,\PP(B_i).
\]
\end{teorema}

\begin{teorema}[Bayes]\label{teo:t4-bayes}
En las condiciones del \cref{teo:t4-total}, si además $\PP(A)>0$,
\[
\PP(B_j\mid A)=\frac{\PP(A\mid B_j)\,\PP(B_j)}{\sum_{i=1}^{n}\PP(A\mid B_i)\,\PP(B_i)}.
\]
\end{teorema}

En la terminología bayesiana, $\PP(B_j)$ es la probabilidad \emph{a priori} de la hipótesis $B_j$, $\PP(A\mid B_j)$ es la \emph{verosimilitud} de la evidencia $A$ bajo esa hipótesis y $\PP(B_j\mid A)$ es la probabilidad \emph{a posteriori}. Con dos hipótesis $B$ y $B^{c}$, dividiendo las dos fórmulas de Bayes se obtiene la \emph{forma de razones} (\emph{odds}):
\[
\frac{\PP(B\mid A)}{\PP(B^{c}\mid A)}=\frac{\PP(B)}{\PP(B^{c})}\cdot\frac{\PP(A\mid B)}{\PP(A\mid B^{c})},
\]
«razón a posteriori = razón a priori $\times$ razón de verosimilitudes».

\addtocounter{definicion}{1}%

\addtocounter{definicion}{1}%

\section{Independencia}

\begin{definicion}[Independencia]
Dos sucesos $A$ y $B$ son \emph{independientes} si $\PP(A\cap B)=\PP(A)\PP(B)$. Si $\PP(B)>0$, esto equivale a $\PP(A\mid B)=\PP(A)$: saber que ha ocurrido $B$ no cambia la probabilidad de $A$. Los sucesos $A_1,\dots,A_n$ son \emph{independientes} (o mutuamente independientes) si para todo subconjunto de índices $S\subseteq[n]$ con $\abs{S}\ge2$,
\[
\PP\Bigl(\bigcap_{i\in S}A_i\Bigr)=\prod_{i\in S}\PP(A_i).
\]
Son \emph{independientes dos a dos} si la condición se cumple al menos para todos los $S$ con $\abs{S}=2$.
\end{definicion}

\begin{proposicion}\label{prop:t4-complementos}
Si $A$ y $B$ son independientes, también lo son $A$ y $B^{c}$, $A^{c}$ y $B$, y $A^{c}$ y $B^{c}$.
\end{proposicion}

\addtocounter{definicion}{1}%

\begin{observacion}[Independencia e incompatibilidad]
No deben confundirse. Si $A$ y $B$ son incompatibles con probabilidades positivas, \emph{no} son independientes: $\PP(A\cap B)=0<\PP(A)\PP(B)$; saber que ocurre $A$ garantiza que no ocurre $B$. La independencia es, además, una propiedad de la ley $\PP$, no de los sucesos como conjuntos: en la práctica suele ser una \emph{hipótesis} de modelización (lanzamientos sucesivos, ensayos repetidos) más que algo que se demuestre~\cite[§17.7.2]{mcs}.
\end{observacion}

Si $\PP(C)>0$, los sucesos $A$ y $B$ son \emph{condicionalmente independientes dado $C$} si $\PP(A\cap B\mid C)=\PP(A\mid C)\PP(B\mid C)$. Ni la independencia implica la independencia condicional ni al revés (\cref{eje:4-ajedrez}).

\addtocounter{definicion}{1}%

\addtocounter{definicion}{1}%


% =====================================================================
\chapter{Variables aleatorias discretas}
\label{tema:5}
% =====================================================================

Una variable aleatoria asigna un número a cada resultado de un experimento; su distribución resume todo lo que interesa de ella. Este tema introduce la función de masa, la esperanza y la varianza, demuestra la linealidad de la esperanza y la explota mediante variables indicadoras, y estudia las distribuciones discretas básicas: Bernoulli, binomial, geométrica, Poisson y categórica. Se siguen los capítulos~3 y~4 de Blitzstein y Hwang~\cite{bh}, el capítulo~2 de Bertsekas y Tsitsiklis~\cite{bt} (clases~5 a~7 de MIT 6.041~\cite{mit6041f10}), el capítulo~4 de Ross~\cite{ross} y el capítulo~18, «Random Variables», de~\cite{mcs}.

\section{Variables aleatorias y función de masa}

\begin{definicion}[Variable aleatoria discreta]
Una \emph{variable aleatoria} es una función $X\colon\Omega\to\R$. Es \emph{discreta} si toma un conjunto finito o numerable de valores. Su \emph{función de masa} es $p_X(x)=\PP(X=x)$, donde $\{X=x\}$ abrevia el suceso $\{\omega\in\Omega: X(\omega)=x\}$, y su \emph{función de distribución} es $F_X(x)=\PP(X\le x)$ para $x\in\R$.
\end{definicion}

La función de masa cumple $p_X(x)\ge0$ y $\sum_xp_X(x)=1$, porque los sucesos $\{X=x\}$ forman una partición de $\Omega$. Para todo conjunto $S\subseteq\R$, $\PP(X\in S)=\sum_{x\in S}p_X(x)$. La función de distribución es no decreciente, tiende a 0 en $-\infty$ y a 1 en $+\infty$, es continua por la derecha y, en el caso discreto, es escalonada, con saltos de altura $p_X(x)$ en cada valor $x$.

\addtocounter{definicion}{1}%

\section{Esperanza y linealidad}

\begin{definicion}[Esperanza]
La \emph{esperanza} (o media) de una variable discreta $X$ es
\[
\E[X]=\sum_{x}x\,p_X(x),
\]
siempre que la serie converja absolutamente. Es el centro de masas de la distribución.
\end{definicion}

\begin{proposicion}[Esperanza de una función: LOTUS]\label{prop:t5-lotus}
Si $X$ es discreta y $g\colon\R\to\R$, entonces $\E[g(X)]=\sum_xg(x)\,p_X(x)$. Si $\Omega$ es discreto, además $\E[X]=\sum_{\omega\in\Omega}X(\omega)\,\PP(\{\omega\})$.
\end{proposicion}

\begin{teorema}[Linealidad de la esperanza]\label{teo:t5-linealidad}
Si $X$ e $Y$ son variables aleatorias discretas definidas sobre el mismo espacio y con esperanza finita, y $a,b\in\R$, entonces
\[
\E[aX+bY]=a\,\E[X]+b\,\E[Y].
\]
No se requiere ninguna hipótesis de independencia.
\end{teorema}

\begin{definicion}[Variable indicadora]
La \emph{indicadora} de un suceso $A$ es la variable $I_A$ que vale 1 si ocurre $A$ y 0 si no.
\end{definicion}

\begin{proposicion}[Puente fundamental]\label{prop:t5-puente}
$\E[I_A]=\PP(A)$. En consecuencia, si $X=I_{A_1}+\dots+I_{A_n}$ cuenta cuántos de los sucesos $A_1,\dots,A_n$ ocurren, $\E[X]=\PP(A_1)+\dots+\PP(A_n)$.
\end{proposicion}

\addtocounter{definicion}{1}%

\section{Varianza e independencia}

\begin{definicion}[Varianza]
La \emph{varianza} de $X$ es $\Var(X)=\E\bigl[(X-\mu)^{2}\bigr]$, con $\mu=\E[X]$, y su \emph{desviación típica} es $\sigma_X=\sqrt{\Var(X)}$. Mide la dispersión de $X$ alrededor de su media, en unidades al cuadrado.
\end{definicion}

\begin{proposicion}\label{prop:t5-var}
$\Var(X)=\E[X^{2}]-\E[X]^{2}$ y, para $a,b\in\R$, $\Var(aX+b)=a^{2}\Var(X)$. En particular $\Var(X)\ge0$ implica $\E[X^{2}]\ge\E[X]^{2}$.
\end{proposicion}

\begin{definicion}[Independencia de variables]
Las variables discretas $X_1,\dots,X_n$ son \emph{independientes} si $\PP(X_1=x_1,\dots,X_n=x_n)=\prod_i\PP(X_i=x_i)$ para todos los valores $x_1,\dots,x_n$.
\end{definicion}

\begin{proposicion}\label{prop:t5-producto}
Si $X$ e $Y$ son discretas, independientes y con esperanza finita, $\E[XY]=\E[X]\E[Y]$. En consecuencia, si $X_1,\dots,X_n$ son independientes (basta dos a dos), $\Var(X_1+\dots+X_n)=\Var(X_1)+\dots+\Var(X_n)$.
\end{proposicion}

\section{Distribuciones discretas básicas}

\begin{definicion}[Bernoulli y binomial]
$X\sim\Ber(p)$, con $p\in[0,1]$, si $\PP(X=1)=p$ y $\PP(X=0)=1-p$: es la indicadora de un «éxito» de probabilidad $p$. $X\sim\Bin(n,p)$ si $X$ es el número de éxitos en $n$ ensayos de Bernoulli independientes con la misma probabilidad $p$.
\end{definicion}

\begin{teorema}\label{teo:t5-binomial}
Si $X\sim\Bin(n,p)$, entonces $\PP(X=k)=\binom nk p^{k}(1-p)^{n-k}$ para $k=0,\dots,n$, $\E[X]=np$ y $\Var(X)=np(1-p)$.
\end{teorema}

\begin{definicion}[Geométrica]
$X\sim\Geom(p)$, con $0<p\le1$, si $X$ es el número de ensayos de Bernoulli independientes de parámetro $p$ necesarios hasta obtener el primer éxito (incluido este). Su función de masa es $\PP(X=k)=(1-p)^{k-1}p$ para $k=1,2,\dots$
\end{definicion}

\begin{proposicion}\label{prop:t5-geom}
Si $X\sim\Geom(p)$: (i) $\PP(X>k)=(1-p)^{k}$; (ii) $\E[X]=1/p$ y $\Var(X)=(1-p)/p^{2}$; (iii) \emph{falta de memoria}: $\PP(X>m+k\mid X>m)=\PP(X>k)$ para enteros $m,k\ge0$.
\end{proposicion}

\begin{definicion}[Poisson]
$X\sim\Poi(\lambda)$, con $\lambda>0$, si $\PP(X=k)=e^{-\lambda}\lambda^{k}/k!$ para $k=0,1,2,\dots$ Se cumple $\E[X]=\Var(X)=\lambda$ (\cref{eje:5-poisson}).
\end{definicion}

La distribución de Poisson modela el número de sucesos raros entre muchas oportunidades independientes: erratas en un libro, llegadas de peticiones a un servidor en un minuto, apariciones de una palabra poco frecuente en un texto largo. El resultado siguiente lo justifica.

\begin{teorema}[Aproximación de Poisson a la binomial]\label{teo:t5-poisson}
Sea $\lambda>0$ y $p_n=\lambda/n$. Para cada $k\ge0$ fijo,
\(
\lim_{n\to\infty}\binom nk p_n^{k}(1-p_n)^{n-k}=e^{-\lambda}\frac{\lambda^{k}}{k!}.
\)
\end{teorema}

\begin{definicion}[Categórica]
Sea $\V=\{v_1,\dots,v_K\}$ un conjunto finito y $\bm p=(p_1,\dots,p_K)$ con $p_i\ge0$ y $\sum_ip_i=1$. Una variable $X$ con valores en $\V$ tiene distribución \emph{categórica}, $X\sim\Cat(\bm p)$, si $\PP(X=v_i)=p_i$. Cuando los valores no son números, en lugar de la esperanza de $X$ se usan las de sus indicadoras: $\E[I_{\{X=v_i\}}]=p_i$.
\end{definicion}

\begin{table}[h]
\centering
\begin{tabular}{@{}lllll@{}}
\toprule
Distribución & Valores & $\PP(X=k)$ & $\E[X]$ & $\Var(X)$ \\
\midrule
$\Ber(p)$ & $\{0,1\}$ & $p^{k}(1-p)^{1-k}$ & $p$ & $p(1-p)$ \\
$\Bin(n,p)$ & $\{0,\dots,n\}$ & $\binom nk p^{k}(1-p)^{n-k}$ & $np$ & $np(1-p)$ \\
$\Geom(p)$ & $\{1,2,\dots\}$ & $(1-p)^{k-1}p$ & $1/p$ & $(1-p)/p^{2}$ \\
$\Poi(\lambda)$ & $\{0,1,\dots\}$ & $e^{-\lambda}\lambda^{k}/k!$ & $\lambda$ & $\lambda$ \\
Uniforme en $[n]$ & $\{1,\dots,n\}$ & $1/n$ & $(n+1)/2$ & $(n^{2}-1)/12$ \\
\bottomrule
\end{tabular}
\caption{Distribuciones discretas básicas.}
\label{tab:t5-discretas}
\end{table}

\addtocounter{definicion}{1}%

\addtocounter{definicion}{1}%


% =====================================================================
\chapter{Variables continuas, desigualdades y ley de los grandes números}
\label{tema:6}
% =====================================================================

Muchas magnitudes, como tiempos, longitudes o pesos de una red neuronal, toman valores en un continuo; su distribución se describe con una densidad en lugar de una función de masa. Este tema introduce las variables continuas y tres distribuciones fundamentales (uniforme, exponencial y normal), y termina con las desigualdades de Markov y Chebyshev y la ley débil de los grandes números, que justifica que una media muestral aproxime una esperanza. Se siguen el capítulo~5 de Blitzstein y Hwang~\cite{bh} (§5.1, §5.2, §5.4 y §5.5) y su §10.1--10.2; las secciones §3.1 a §3.3 y §5.1 a §5.2 de Bertsekas y Tsitsiklis~\cite{bt} (clases~8 y~19 de MIT 6.041~\cite{mit6041f10}); los capítulos~5 y~8 de Ross~\cite{ross}; el capítulo~19 de~\cite{mcs} y el capítulo~8 de Grinstead y Snell~\cite{gs}.

\section{Densidad y función de distribución}

\begin{definicion}[Variable continua]
Una variable aleatoria $X$ es \emph{continua} si existe una función $f_X\colon\R\to[0,\infty)$, su \emph{densidad}, tal que para todos $a\le b$
\[
\PP(a\le X\le b)=\int_a^bf_X(x)\,dx.
\]
Necesariamente $\int_{-\infty}^{\infty}f_X(x)\,dx=1$. La \emph{función de distribución} es $F_X(x)=\PP(X\le x)=\int_{-\infty}^{x}f_X(t)\,dt$.
\end{definicion}

Consecuencias inmediatas: $\PP(X=x)=\int_x^xf_X=0$ para todo $x$, de modo que da igual escribir $<$ o $\le$; $F_X$ es continua y no decreciente, y en los puntos donde $f_X$ es continua, $F_X'(x)=f_X(x)$ por el teorema fundamental del cálculo. La densidad \emph{no} es una probabilidad y puede ser mayor que~1: lo que se interpreta es $\PP(x\le X\le x+\delta)\approx f_X(x)\,\delta$ para $\delta$ pequeño.

\begin{definicion}[Esperanza y varianza]
Para $X$ continua, $\E[X]=\int_{-\infty}^{\infty}x\,f_X(x)\,dx$ (si la integral converge absolutamente) y, para una función $g$, $\E[g(X)]=\int g(x)f_X(x)\,dx$. La varianza se define como en el caso discreto, $\Var(X)=\E[(X-\E X)^{2}]=\E[X^{2}]-\E[X]^{2}$, y la linealidad y las propiedades de la \cref{prop:t5-var} siguen siendo válidas.
\end{definicion}

\section{Distribuciones uniforme y exponencial}

\begin{definicion}[Uniforme]
$X\sim\Unif(a,b)$, con $a<b$, si $f_X(x)=\frac{1}{b-a}$ en $[a,b]$ y $0$ fuera. Entonces $F_X(x)=\frac{x-a}{b-a}$ para $x\in[a,b]$: la probabilidad de un subintervalo es proporcional a su longitud.
\end{definicion}

\begin{proposicion}\label{prop:t6-unif}
Si $X\sim\Unif(a,b)$, entonces $\E[X]=\frac{a+b}{2}$ y $\Var(X)=\frac{(b-a)^{2}}{12}$.
\end{proposicion}

\begin{definicion}[Exponencial]
$X\sim\Expo(\lambda)$, con $\lambda>0$, si $f_X(x)=\lambda e^{-\lambda x}$ para $x\ge0$ y $0$ para $x<0$. Integrando, $F_X(x)=1-e^{-\lambda x}$ y $\PP(X>x)=e^{-\lambda x}$ para $x\ge0$. Se cumple $\E[X]=1/\lambda$ y $\Var(X)=1/\lambda^{2}$ (\cref{eje:6-expo}).
\end{definicion}

\begin{proposicion}[Falta de memoria]\label{prop:t6-memoria}
Si $X\sim\Expo(\lambda)$, entonces $\PP(X>s+t\mid X>s)=\PP(X>t)$ para todos $s,t\ge0$.
\end{proposicion}

La exponencial es la versión continua de la geométrica: modela tiempos de espera hasta un suceso que «no envejece», como la llegada de la siguiente petición a un servidor.

\section{Distribución normal y tipificación}

\begin{definicion}[Normal]
$Z\sim\Nor(0,1)$ (\emph{normal estándar}) si su densidad es $\varphi(z)=\frac{1}{\sqrt{2\pi}}e^{-z^{2}/2}$; su función de distribución se denota $\Phi$. Para $\mu\in\R$ y $\sigma>0$, $X\sim\Nor(\mu,\sigma^{2})$ si
\[
f_X(x)=\frac{1}{\sigma\sqrt{2\pi}}\exp\Bigl(-\frac{(x-\mu)^{2}}{2\sigma^{2}}\Bigr).
\]
\end{definicion}

Que $\varphi$ integra~1 equivale a la integral gaussiana $\int_{-\infty}^{\infty}e^{-z^{2}/2}dz=\sqrt{2\pi}$, que se demuestra en el módulo~5 con coordenadas polares. Por la simetría de $\varphi$, $\E[Z]=0$ y $\Phi(-z)=1-\Phi(z)$; integrando por partes se obtiene $\E[Z^{2}]=1$, luego $\Var(Z)=1$. $\Phi$ no tiene expresión con funciones elementales y se tabula (\cref{tab:t6-phi}).

\begin{teorema}[Tipificación]\label{teo:t6-tipificacion}
Si $Z\sim\Nor(0,1)$, entonces $X=\mu+\sigma Z\sim\Nor(\mu,\sigma^{2})$. Recíprocamente, si $X\sim\Nor(\mu,\sigma^{2})$, la variable tipificada $Z=(X-\mu)/\sigma$ es $\Nor(0,1)$ y
\[
\PP(X\le x)=\Phi\Bigl(\frac{x-\mu}{\sigma}\Bigr).
\]
En particular, $\E[X]=\mu$ y $\Var(X)=\sigma^{2}$.
\end{teorema}

\begin{table}[h]
\centering
\begin{tabular}{@{}lcccccccc@{}}
\toprule
$z$ & $0{,}5$ & $1$ & $1{,}5$ & $1{,}645$ & $1{,}96$ & $2$ & $2{,}5$ & $3$ \\
\midrule
$\Phi(z)$ & $0{,}6915$ & $0{,}8413$ & $0{,}9332$ & $0{,}9500$ & $0{,}9750$ & $0{,}9772$ & $0{,}9938$ & $0{,}9987$ \\
\bottomrule
\end{tabular}
\caption{Valores de la función de distribución normal estándar. Para $z<0$, $\Phi(z)=1-\Phi(-z)$.}
\label{tab:t6-phi}
\end{table}

\addtocounter{definicion}{1}%

\section{Desigualdades de Markov y Chebyshev}

\begin{teorema}[Desigualdad de Markov]\label{teo:t6-markov}
Si $X\ge0$ y $a>0$, entonces $\PP(X\ge a)\le\dfrac{\E[X]}{a}$.
\end{teorema}

\begin{teorema}[Desigualdad de Chebyshev]\label{teo:t6-chebyshev}
Si $X$ tiene media $\mu$ y varianza $\sigma^{2}$ finitas, para todo $c>0$,
\[
\PP(\abs{X-\mu}\ge c)\le\frac{\sigma^{2}}{c^{2}}.
\]
\end{teorema}

Con $c=k\sigma$, Chebyshev dice que $\PP(\abs{X-\mu}\ge k\sigma)\le1/k^{2}$ para \emph{cualquier} distribución: a dos desviaciones típicas hay como mucho un $25\,\%$ de la probabilidad. Para la normal el valor exacto es un $4{,}6\,\%$; la cota es peor porque no usa nada de la forma de la distribución.

\section{Ley débil de los grandes números}

\begin{proposicion}[Media muestral]\label{prop:t6-media}
Sean $X_1,\dots,X_n$ independientes con la misma media $\mu$ y la misma varianza $\sigma^{2}$, y sea $\Xbar_n=\frac1n\sum_{i=1}^{n}X_i$. Entonces $\E[\Xbar_n]=\mu$ y $\Var(\Xbar_n)=\sigma^{2}/n$.
\end{proposicion}

\begin{teorema}[Ley débil de los grandes números]\label{teo:t6-lgn}
En las condiciones de la \cref{prop:t6-media}, para todo $\varepsilon>0$,
\[
\PP\bigl(\abs{\Xbar_n-\mu}\ge\varepsilon\bigr)\le\frac{\sigma^{2}}{n\varepsilon^{2}}\xrightarrow[n\to\infty]{}0.
\]
Se dice que $\Xbar_n$ \emph{converge en probabilidad} a $\mu$.
\end{teorema}

\begin{corolario}[Frecuencias relativas]\label{cor:t6-frecuencias}
Si $A$ es un suceso de probabilidad $p$ y $f_n$ es la frecuencia relativa con que ocurre en $n$ repeticiones independientes, entonces $\PP(\abs{f_n-p}\ge\varepsilon)\le\dfrac{p(1-p)}{n\varepsilon^{2}}\le\dfrac{1}{4n\varepsilon^{2}}$.
\end{corolario}

El corolario reconcilia la definición axiomática con la interpretación frecuentista: la probabilidad de un suceso es el valor al que tiende, en probabilidad, su frecuencia relativa. Con $\varepsilon=0{,}01$ y $n=10^{5}$, la cota es $1/(4\cdot10^{5}\cdot10^{-4})=0{,}025$.

\addtocounter{definicion}{1}%

\addtocounter{definicion}{1}%



\backmatter
% =====================================================================
\begin{thebibliography}{99}
\addcontentsline{toc}{chapter}{Bibliografía}
\markboth{Bibliografía}{}
\small

\bibitem{bt}
D.~P. Bertsekas y J.~N. Tsitsiklis.
\newblock \emph{Introduction to Probability}, 2.ª ed.
\newblock Athena Scientific, 2008.
\newblock Texto del curso 6.041 del MIT.

\bibitem{bh}
J.~K. Blitzstein y J.~Hwang.
\newblock \emph{Introduction to Probability}, 2.ª ed.
\newblock Chapman \& Hall/CRC Press, 2019.
\newblock Texto del curso Stat~110 de Harvard.

\bibitem{stat110}
J.~K. Blitzstein.
\newblock Stat 110: Probability.
\newblock Harvard University. \emph{Strategic Practice and Homework}~1 a~6 y~10, otoño de 2011.
\newblock \url{https://stat110.hsites.harvard.edu/strategic-practice-problems}

\bibitem{chen}
M.~Chen \emph{et al.}
\newblock Evaluating large language models trained on code.
\newblock arXiv:2107.03374, 2021.
\newblock \url{https://arxiv.org/abs/2107.03374}

\bibitem{gs}
C.~M. Grinstead y J.~L. Snell.
\newblock \emph{Introduction to Probability}, 2.ª ed. revisada.
\newblock American Mathematical Society. Versión libre del proyecto CHANCE de 4 de julio de 2006, bajo licencia GNU FDL.
\newblock \url{https://chance.dartmouth.edu/teaching_aids/books_articles/probability_book/book.html}

\bibitem{mcs}
E.~Lehman, F.~T. Leighton y A.~R. Meyer.
\newblock \emph{Mathematics for Computer Science}.
\newblock Edición del curso 6.042J del MIT, primavera de 2015 (MIT OpenCourseWare). Se citan los capítulos~14 y~16 a~19 de esa edición.
\newblock \url{https://ocw.mit.edu/courses/6-042j-mathematics-for-computer-science-spring-2015/pages/readings/}

\bibitem{mit6042}
A.~R. Meyer y A.~Chlipala.
\newblock 6.042J: Mathematics for Computer Science.
\newblock Massachusetts Institute of Technology, MIT OpenCourseWare, primavera de 2015.
\newblock \url{https://ocw.mit.edu/courses/6-042j-mathematics-for-computer-science-spring-2015/}

\bibitem{mit1805}
J.~Orloff y J.~Bloom.
\newblock 18.05: Introduction to Probability and Statistics.
\newblock Massachusetts Institute of Technology, MIT OpenCourseWare, primavera de 2022. \emph{Problem Sets}~1 a~3.
\newblock \url{https://ocw.mit.edu/courses/18-05-introduction-to-probability-and-statistics-spring-2022/}

\bibitem{ross}
S.~Ross.
\newblock \emph{A First Course in Probability}, 10.ª ed.
\newblock Pearson, 2024.

\bibitem{mit6041f10}
J.~N. Tsitsiklis.
\newblock 6.041: Probabilistic Systems Analysis and Applied Probability.
\newblock Massachusetts Institute of Technology, MIT OpenCourseWare, otoño de 2010. \emph{Problem Sets}~1 a~4 y~9.
\newblock \url{https://ocw.mit.edu/courses/6-041-probabilistic-systems-analysis-and-applied-probability-fall-2010/}

\end{thebibliography}


\end{document}
